\section*{14. 环形房间化学品存放问题}

给定$n$个首尾相连的环形房间，每个房间容量为$w_i$（杂物房间容量为0）。约束：化学品不能存放在相邻房间。求最多能存放的化学品总容量。

\textbf{线性序列DP（打家劫舍）}：
\begin{itemize}
    \item \textbf{状态定义}：$dp[i]$表示前$i$个房间的最大容量
    \item \textbf{转移方程}：$dp[i] = \max(dp[i-1], dp[i-2] + w_i)$
    \item \textbf{边界条件}：$dp[0] = 0$，$dp[1] = w_1$
\end{itemize}

\textbf{环形转线性}：首尾相邻，拆分为两个线性问题（不选第1个或不选第n个）

\begin{lstlisting}[language=C++, style=cppstyle]
// 线性序列DP（打家劫舍）
int houseRobberLinear(vector<int>& w, int start, int end) {
    int n = end - start + 1;
    if (n == 0) return 0;
    if (n == 1) return w[start];

    int prev2 = 0;      // dp[i-2]
    int prev1 = w[start];  // dp[i-1]

    for (int i = start + 1; i <= end; i++) {
        int curr = max(prev1, prev2 + w[i]);
        prev2 = prev1;
        prev1 = curr;
    }

    return prev1;
}

// 环形房间问题主函数
int houseRobberCircular(vector<int>& w) {
    int n = w.size();
    if (n == 0) return 0;
    if (n == 1) return w[0];
    if (n == 2) return max(w[0], w[1]);

    // 情况1：不选第1个房间，考虑2~n
    int case1 = houseRobberLinear(w, 1, n - 1);

    // 情况2：不选第n个房间，考虑1~n-1
    int case2 = houseRobberLinear(w, 0, n - 2);

    return max(case1, case2);
}
\end{lstlisting}


\begin{enumerate}
    \item 环形拆成两个线性问题：不选第1个、不选第n个
    \item 线性DP：dp[i]表示前i个房间的最大容量
    \item 比较"不选i"和"选i"
    \item 取两种情况的最大值
\end{enumerate}

